\scanpage{16}

\mspage{8)}

\diagram{\begin{tikzcd}[ampersand replacement=\&,column sep=small,row sep=small]
x\vee x \arrow[r,tail,"{\deleted{\unsure{x|f|(}} f_0\vee f'_0\deleted{\unsure{g|)}}}"] \arrow[d,"{\nabla_x}"'] \arrow[dr,phantom,"{\mathrm{cocart}}"{description}] \& \bar y\amalg\bar y' \arrow[d,dashed] \arrow[drr,"{(\alpha,\alpha')}"] \& {} \& {} \\
x \arrow[r,tail,dashed,"g"] \arrow[rrr,bend right=17,"f"'] \& \widetilde y \arrow[r,dashed,"{\varphi_0}"] \arrow[rr,dashed,"{\deleted{\unsure{\varphi|\alpha}}\varphi}"'] \& \bar y'' \arrow[r,dashed,"{\alpha''}","{\sim}"'] \& y
\end{tikzcd}}


on factorise $\varphi$ en $\deleted{\unsure{\varphi|\alpha}}\alpha''\varphi_0$ avec $\varphi_0\in\mathrm{cof}$ $\deleted{\unsure{\varphi|\alpha}}\alpha''\in W$, et on prend
\[
 f_0''=\varphi_0 g
\]

Mais si $\alpha,\alpha'$ ont des sections (il suffirait que $\alpha$ ou $\alpha'$ les \unsure{aient|ait}) \deleted{\unsure{il reste donc|On montre donc|il reste une}} il est trivial que \deleted{\unsure{les autres le|il en est de|les autres la}} $\alpha''$ (quel lui \unsure{qu'il soit|aussi soit}) en a aussi une. Donc on est ramené au cas où on a une flèche $u:(f_0,\alpha)\longrightarrow(\unsure{f'_0|f_0|f''_0},\alpha')$ dans la catégorie des factorisations de $f$ en $\alpha f_0$. Donc \uline{Si} on choisit la section $\unsure{i^*|i|i'}$ pour $\alpha^*$ comme l'image par $u$ de la section $i'$ de $\alpha'$ \deleted{\unsure{par u|on}}, on gagne.

\diagram{\begin{tikzcd}[ampersand replacement=\&,column sep=small,row sep=small]
a \arrow[rr,"f"] \arrow[dr,"{f'_0}"'] \arrow[drr,"{f^*_0}"] \& {} \& y \\
{} \& \bar y' \arrow[ur,"{\alpha'}"] \arrow[r,leftrightarrow,"u"'] \& \bar y^* \arrow[u,"{\alpha^*}"]
\end{tikzcd}}

Donc il reste à prouver que \deleted{le choix de} pour une factorisation choisie de $f$ en $\alpha f_0$ comme dessus, $[f_0,i]$ ne dépend pas du choix de la section $i$. Comme $[f_0,i]=[\mathrm{id}_{\bar y},i]\circ(f,\mathrm{id}_{\bar y}]$, On est ramené à ceci :

\diagram{\begin{tikzcd}[ampersand replacement=\&,column sep=small,row sep=small]
x \arrow[rr,"{\unsure{f^{\sim}|f^{\nearrow}}}"] \arrow[dr,"{f_0}"'] \& {} \& y \arrow[dl,bend left=18,"i"] \arrow[dl,bend left=40,"{i'}"] \\
{} \& \bar y \arrow[ur,"{\alpha}"',"{\sim}"] \& {}
\end{tikzcd}}

$(\unsure{\ast_2|\ast\ast|\ast})$ Soit $x\xrightarrow[\sim]{f}y$ de $W$, tel que $f$ ait une \deleted{\unsure{(deux $i,i'$ de $W$)|(deux $i,i'(\ldots)$)|(deux $i,i'\in W$)}} section $i$, alors \inserted{la classe dans}\deleted{\unsure{pour deux telles sections $i,i'$|pour deux telles sections $i,i'$)}} \deleted{\unsure{les comp. connexes de|les composantes connexes de}} $\pi_0(\underline{\mathrm{Ch}}_0(x,y))$ \deleted{de} de $(1_x,i)$

\diagram{\begin{tikzcd}[ampersand replacement=\&,column sep=small,row sep=small]
\deleted{\unsure{x|y}}\quad x \arrow[dr,"{1_x}"'] \& y \arrow[d,"i"] \\
{} \& x
\end{tikzcd}}

ne dépend pas du choix de la section $i$.

Cette condition est (néc. et) suffis$^{\mathrm{te}}$ pour que $\widetilde C\longrightarrow W^{-1}C$ soit fidèle (donc un iso), comme on doit pouvoir vérifier aisément, en construisant bel et bien le foncteur cherché $C\longrightarrow\widetilde C$.

% archive sheet number: 15
